\hebrewsection{Audio and Speech Processing}

Audio and speech signals are one-dimensional signals that carry a wealth of information. This chapter explores the production of speech and the algorithms used for its analysis and synthesis.

\subsection{Digital Audio Foundations}
High-fidelity audio requires high sampling rates (e.g., 44.1 kHz or 48 kHz) and high quantization depth (16 to 24 bits).

\subsection{Speech Production Model}
The human speech production system can be modeled as an excitation source (vocal folds or turbulent airflow) driving a time-varying filter (the vocal tract).

\subsubsection{Linear Predictive Coding (LPC)}
LPC is a powerful method for speech analysis. It assumes the vocal tract is an all-pole filter:
\begin{equation}
    s[n] = \sum_{i=1}^{p} a_i s[n-i] + e[n]
\end{equation}
By estimating the coefficients $a_i$, we can characterize the resonances (formants) of the vocal tract. LPC is used in low-bitrate speech codecs like those in early cellular systems.

\subsection{Speech Analysis: The Mel-Frequency Cepstrum}
The human ear perceives frequency non-linearly. The \textbf{Mel Scale} approximates this perception.
\subsubsection{Mel-Frequency Cepstral Coefficients (MFCCs)}
MFCCs are the standard features for speech and audio recognition. They are calculated by:
1. Taking the STFT.
2. Mapping the power to the Mel scale.
3. Taking the log of the powers.
4. Taking the Discrete Cosine Transform of the log powers.

\subsection{Audio Compression (Perceptual Coding)}
MP3 and AAC use \textbf{psychoacoustic models} to remove sounds that are masked by other louder sounds or are outside the human hearing range.
\begin{itemize}
    \item \textbf{Frequency Masking:} A loud tone at 1 kHz masks quieter tones at 1.1 kHz.
    \item \textbf{Temporal Masking:} A loud sound masks quieter sounds that occur shortly before or after it.
\end{itemize}

\subsection{Speech Recognition: From HMM to Deep Learning}
Traditionally, speech recognition used \textbf{Hidden Markov Models (HMMs)} to model the temporal sequence of phonemes. Modern systems use \textbf{Recurrent Neural Networks (RNNs)} and \textbf{Transformers} (like OpenAI's Whisper) to achieve human-level accuracy.

\subsection{Music Information Retrieval (MIR)}
MIR involves extracting high-level information from music audio, such as:
\begin{itemize}
    \item \textbf{Pitch Detection:} Identifying the fundamental frequency.
    \item \textbf{Beat Tracking:} Identifying the tempo and pulse.
    \item \textbf{Genre Classification:} Using machine learning on spectral features.
\end{itemize}
